% Sección: Ejemplos de uso de teoremas, definiciones y corolarios
\section{Teoremas y Definiciones Fundamentales}\label{sec:teoremas}

\begin{definition}[Coseno de un ángulo]
Sea $\theta$ un ángulo en el círculo unitario. El \textbf{coseno} de $\theta$ se define como la coordenada $x$ del punto donde el lado terminal del ángulo interseca el círculo unitario.
\end{definition}

\begin{definition}[Seno de un ángulo]
Sea $\theta$ un ángulo en el círculo unitario. El \textbf{seno} de $\theta$ se define como la coordenada $y$ del punto donde el lado terminal del ángulo interseca el círculo unitario.
\end{definition}

\begin{theorem}[Identidad Fundamental del Coseno de Suma]\label{thm:coseno_suma}
Para cualquier par de ángulos $a$ y $b$:
\begin{equation}\label{eq:teorema_coseno_suma}
\cos(a + b) = \cos a \cos b - \sin a \sin b
\end{equation}
\end{theorem}

\begin{proof}
Esta identidad puede demostrarse de múltiples formas:
\begin{enumerate}
\item Usando la multiplicación de matrices de rotación (ver sección \ref{sec:demcos})
\item Mediante geometría del círculo unitario
\item Utilizando la fórmula de Euler y números complejos
\end{enumerate}
Las demostraciones detalladas se encuentran en las secciones anteriores.
\end{proof}

\begin{corollary}[Identidad del Coseno de Diferencia]
Para cualquier par de ángulos $a$ y $b$:
\begin{equation}\label{eq:coseno_diferencia}
\cos(a - b) = \cos a \cos b + \sin a \sin b
\end{equation}
\end{corollary}

\begin{proof}
Aplicando el Teorema \ref{thm:coseno_suma} con $b$ reemplazado por $-b$:
\begin{align}
\cos(a - b) &= \cos(a + (-b)) \\
&= \cos a \cos(-b) - \sin a \sin(-b) \\
&= \cos a \cos b - \sin a (-\sin b) \\
&= \cos a \cos b + \sin a \sin b
\end{align}
donde usamos que $\cos(-b) = \cos b$ y $\sin(-b) = -\sin b$.
\end{proof}

\begin{lemma}[Propiedades de Paridad]\label{lemma:paridad}
Para cualquier ángulo $\theta$:
\begin{align}
\cos(-\theta) &= \cos \theta \quad \text{(función par)} \\
\sin(-\theta) &= -\sin \theta \quad \text{(función impar)}
\end{align}
\end{lemma}

\begin{proposition}[Matriz de Rotación]
La rotación de un punto en el plano por un ángulo $\theta$ puede expresarse mediante la matriz:
\begin{equation}\label{eq:matriz_rotacion_prop}
R(\theta) = \begin{pmatrix}
\cos \theta & -\sin \theta \\
\sin \theta & \cos \theta
\end{pmatrix}
\end{equation}
\end{proposition}

\begin{example}[Aplicación de la Identidad]
Calculemos $\cos(75^\circ)$ usando $\cos(75^\circ) = \cos(45^\circ + 30^\circ)$:
\begin{align}
\cos(75^\circ) &= \cos(45^\circ + 30^\circ) \\
&= \cos 45^\circ \cos 30^\circ - \sin 45^\circ \sin 30^\circ \\
&= \frac{\sqrt{2}}{2} \cdot \frac{\sqrt{3}}{2} - \frac{\sqrt{2}}{2} \cdot \frac{1}{2} \\
&= \frac{\sqrt{6}}{4} - \frac{\sqrt{2}}{4} \\
&= \frac{\sqrt{6} - \sqrt{2}}{4}
\end{align}
\end{example}

\begin{remark}
La identidad del coseno de suma es fundamental en trigonometría y tiene aplicaciones en:
\begin{itemize}
\item Análisis de Fourier
\item Procesamiento de señales
\item Mecánica cuántica
\item Cristalografía
\end{itemize}
\end{remark}

\begin{note}
El paquete \texttt{amsthm} permite definir diferentes estilos:
\begin{itemize}
\item \textbf{plain}: Para teoremas, lemas, corolarios (texto en cursiva)
\item \textbf{definition}: Para definiciones y ejemplos (texto normal)
\item \textbf{remark}: Para observaciones y notas (texto normal)
\end{itemize}
\end{note}
